
\section{\BookRecordIndex{T0186}交点正性}\label{sec:E-2}

现在从 \BookRecordIndex{T0036}Micallef--White 定理\BookRecordIndex{T0326}推出定理~\ref{thm:E-1-5} 和
定理~\ref{thm:E-1-6}。出发点是\BookRecordIndex{T0201}局部交数的定义。

\begin{definition}\label{def:E-2-1}
设 $M$ 为定向四维流形，$\Sigma_0,\Sigma_1$ 为两个定向二维流形，
$u_0:\Sigma_0\to M$、$u_1:\Sigma_1\to M$ 为光滑映射。
若 $(z_0,z_1)\in\Sigma_0\times\Sigma_1$ 满足
$u_0(z_0)=u_1(z_1)$，且 $z_i$ 分别有闭包紧的开邻域 $U_i\subset\Sigma_i$，
使对 $w_i\in\overline U_i$，
\[
  u_0(w_0)=u_1(w_1)\quad\Longrightarrow\quad w_0=z_0,\ w_1=z_1,
\]
则称 $(z_0,z_1)$ 为 $u_0,u_1$ 的一个\emph{孤立交点}。
在孤立交点 $(z_0,z_1)$ 处，两映射的\emph{\BookRecordIndex{T0201}局部交数}定义为
\[
  \iota(u_0,u_1;z_0,z_1):=v_0\cdot v_1,
\]
其中 $W\subset M$ 是 $u_0(z_0)=u_1(z_1)$ 的紧致可缩邻域，并满足
\[
  W\cap u_0(\partial U_0)=W\cap u_1(\partial U_1)=\varnothing;
\]
而 $v_i:U_i\to M$（$i=0,1$）是相互横截的光滑映射，且
$v_i$ 在 $U_i\setminus u_i^{-1}(W)$ 上与 $u_i$ 一致，并满足
$v_i(U_i\cap u_i^{-1}(W))\subset W$。

因为 $W$ 可缩，$v_0$ 的任意两种选取可通过一个取值于 $W$、支集包含在
$U_0\cap u_0^{-1}(W)$ 中的同伦相连。因此，\BookRecordIndex{T0201}局部交数与定义它所用的扰动
$v_0,v_1$ 的选择无关；随之也与邻域 $U_0,U_1,W$ 的选择无关。注意，
定义~\ref{def:E-2-1} 可原样推广到光滑映射 $u:\Sigma\to M$ 的自交点
$(z_0,z_1)\in\Sigma\times\Sigma$，其中 $u(z_0)=u(z_1)$ 且 $z_0\ne z_1$
（分别用 $z_0,z_1$ 的不交邻域代替 $\Sigma_0,\Sigma_1$）。
\end{definition}

\begin{proposition}\label{prop:E-2-2}
设 $(M,J)$ 为四维殆复流形，$(\Sigma,j)$ 为黎曼曲面，不要求闭或连通，
$u:\Sigma\to M$ 为单纯 \BookRecordIndex{T0003}$J$-全纯曲线。则 $u$ 的每个自交点
$(z_0,z_1)\in\Sigma\times\Sigma\setminus\Delta$ 都是孤立的，且满足
$\iota(u,u;z_0,z_1)\ge1$；等号成立当且仅当该交点横截。
\end{proposition}

为证明此命题，我们构造 $u$ 的适当扰动。形如
\[
  u(z)=z^k(1,p(z))
\]
的多项式 $u:\mathbb C\to\mathbb C\times\mathbb C^{n-1}$，其中 $p(0)=0$，
若在零的某个邻域上的限制为单射，就称为\emph{单纯的}。
等价地，对每个不等于 $1$ 的 $k$ 次单位根 $\zeta$，
都有 $p(z)\not\equiv p(\zeta z)$\sourcefootnote{原书此处写作“对每个 $k$ 次单位根 $\zeta$，$p(z)\not\equiv p(\zeta z)$”，漏排 $\zeta\ne1$；取 $\zeta=1$ 时该条件不可能成立。据前一定义及定理 E.1.2 的证明补上排除条件。}
（见定理~\ref{thm:E-1-2} 的证明）。又等价地，$u$ 不能写成
$u(z)=v(z^m)$，其中 $v$ 为多项式，$m>1$ 为整除 $k$ 的整数。

\begin{lemma}\label{lem:E-2-3}
设 $u:\mathbb D\to\mathbb C\times\mathbb C^{n-1}$ 为形如
$u(z)=z^k(1,p(z))$ 的单纯多项式，其中 $k$ 为正整数且 $p(0)=0$。
设 $U\subset\mathbb C$ 是零的开邻域，$V\subset\mathbb C$、
$W\subset\mathbb C^n$ 是零的闭邻域，满足
\[
  V\subset U,\qquad u(V)\subset\operatorname{int}(W),\qquad
  W\cap u(\partial U)=\varnothing.
\]
则对每个 $\varepsilon>0$，存在浸入 $v:U\to\mathbb C^n$，满足：
\begin{enumerate}
  \item[(i)] $v$ 在 $U\setminus u^{-1}(W)$ 上与 $u$ 一致，且
    $v(U\cap u^{-1}(W))\subset W$；
  \item[(ii)] $v$ 在 $V$ 中全纯，且 $\|u-v\|_{C^1}<\varepsilon$；
  \item[(iii)] $\#\{v^{-1}(0)\}=k$。
\end{enumerate}
\end{lemma}
\begin{proof}
定义
\begin{equation}\label{eq:E-2-1}
  u_\varepsilon(z):=f_\varepsilon(z)(1,p(z)),\qquad
  f_\varepsilon(z):=\prod_{j=0}^{k-1}(z+j\varepsilon).
\end{equation}
取零的一个开邻域 $U_1\subset\mathbb C$，使
\[
  z\in U_1\setminus\{0\},\quad \lambda^k=1,\quad\lambda\ne1
  \quad\Longrightarrow\quad p'(z)\ne0,\quad p(z)\ne p(\lambda z).
\]
由于 $f'_\varepsilon(0)\ne0$，$f_\varepsilon$ 与
$f'_\varepsilon$ 没有公共零点，而且
\[
  u'_\varepsilon=(f'_\varepsilon,f'_\varepsilon p+f_\varepsilon p'),
\]
所以 $u_\varepsilon$ 在 $U_1$ 上的限制是浸入。此外，对每个
$\varepsilon\ne0$，有\sourcefootnote{原书在 $f_\varepsilon(z)=\prod_{j=0}^{k-1}(z+j\varepsilon)$ 之后将零点印作 $\{j\varepsilon\}$；由乘积中的正号可唯一确定为 $\{-j\varepsilon\}$，故据式 (E.2.1) 改正。}
\[
  u_\varepsilon^{-1}(0)=\{-j\varepsilon\mid j=0,\ldots,k-1\}.
\]

现在把 $u_\varepsilon$ 与 $u$ 拼接起来，得到 $v$。取开集
$V'\subset U$，使 $V\subset V'$ 且
$\overline{V'}\subset\operatorname{int}(u^{-1}(W))$。
选取光滑\BookRecordIndex{T0188}截断函数 $\beta:U\to[0,1]$，满足
\[
  \beta(z)=\begin{cases}
    1,&z\in V,\\
    0,&z\in U\setminus V'.
  \end{cases}
\]
定义
\begin{equation}\label{eq:E-2-2}
  v_\varepsilon(z):=u(z)+\beta(z)(u_\varepsilon(z)-u(z)).
\end{equation}
当 $\varepsilon>0$ 足够小时，映射 $v:=v_\varepsilon:U\to\mathbb C^n$\sourcefootnote{原书此处把 $v$ 的靶印作 $\mathbb C^2$，但本引理设 $u:U\to\mathbb C^n$，式 (E.2.2) 也在 $\mathbb C^n$ 中定义 $v_\varepsilon$，故据引理陈述改作 $\mathbb C^n$。}
是浸入，在 $U\setminus u^{-1}(W)$ 上与 $u$ 一致，且满足
$v(U\cap u^{-1}(W))\subset W$。\sourcefootnote{原书证明末句将上述两个定义域写作 $U_1\setminus u^{-1}(W)$ 与 $U_1\cap u^{-1}(W)$，而引理 (i) 要求在 $U$ 上成立；由于\BookRecordIndex{T0188}截断函数定义在 $U$ 且支集位于 $V'\subset u^{-1}(W)$，据构造改作 $U$。}
由构造可知，它也满足 (ii)、(iii)：除 $\beta$ 的支集内可能不是
$J$-全纯的，在该支集中的点处，$v$ 可任意地在 $C^1$ 意义下接近 $u$。
\end{proof}

\begin{proof}[命题~\ref{prop:E-2-2} 的证明]
对 $i=0,1$，用 $u_i:\Sigma_i\to M$ 表示 $u$ 在 $z_i$ 的一个邻域上的限制。
由定理~\ref{thm:E-1-2} 中的局部论证\sourcefootnote{原书此处直接引用定理 E.1.2，但本命题允许 $\Sigma$ 非闭，而定理 E.1.2 的陈述要求 $\Sigma$ 闭；其证明中关于自交点孤立性的论证是局部的，不用闭性。此处据该证明改称“局部论证”。}，
$(z_0,z_1)$ 是 $u_0,u_1$ 的孤立交点。
由定理~\ref{thm:E-1-1}，不失一般性可令
$M=\mathbb C^2$、$\Sigma_0=\Sigma_1=\mathbb C$、$z_0=z_1=0$，
且存在矩阵 $L_i$，使 $L_i u_i$ 是形如
\[
  L_i u_i(z)=z^{k_i}(1,p_i(z))
\]
的多项式，其中 $k_i$ 为正整数且 $p_i(0)=0$。

\noindent\textbf{第一步：$k_0=k_1=1$ 时结论成立。}
若两曲线横截相交，结论显然。否则，由定理~\ref{thm:E-1-1} 可设
$L_0=L_1=\mathbb 1$。以
$\Psi(z,w)=(z,w-zp_0(z))$ 将 $u_i$ 换成 $\Psi\circ u_i$，
还可设 $u_0(z)=(z,0)$。于是
\[
  u_0(z)=(z,0),\qquad u_1(z)=(z,z^kq(z)),
\]
其中 $k\ge2$，$q:\mathbb C\to\mathbb C$ 为多项式且 $q(0)\ne0$。
$u_1$ 的第二分量 $z^kq$ 是一个在 $z=0$ 处有 $k$ 重零点的多项式。
因此，只要把它扰动成在圆盘 $|z|\le\varepsilon$ 中有 $k$ 个互异零点，
便得到与 $u_0$ 相交 $k$ 次的映射 $v_1$。具体如下：取零的一个开邻域
$U_1\subset\mathbb C$，使对 $z\in\overline U_1$ 有 $q(z)\ne0$；
令 $\beta:\mathbb C\to[0,1]$ 为支集包含在 $U_1$ 中、在零附近等于 $1$
的光滑\BookRecordIndex{T0188}截断函数，$f_\varepsilon$ 如~\eqref{eq:E-2-1}，并定义
\[
  v_1(z):=\bigl(z,\bigl(z^k+\beta(z)(f_\varepsilon(z)-z^k)\bigr)q(z)\bigr).
\]
当 $\varepsilon>0$ 足够小时，$v_1$ 是浸入，且
$\iota(u_0,u_1;0,0)=u_0\cdot v_1=k\ge2$，如所需。

\noindent\textbf{第二步：$k_0k_1>1$ 时结论成立。}
由假设，存在以零为中心的开圆盘 $U\subset\mathbb C$，使对
$w_0,w_1\in U\setminus\{0\}$ 有 $u_0(w_0)\ne u_1(w_1)$。
取以零为中心的闭球 $W\subset\mathbb C^2$，使
$W\cap u_0(\partial U)=W\cap u_1(\partial U)=\varnothing$。
再取闭圆盘 $V_0,V_1\subset U$，使
$u_0(V_0),u_1(V_1)\subset\operatorname{int}(W)$。
选取 $\varepsilon>0$，使
\[
  \inf_{\substack{w_0\in U\\w_1\in U\setminus V_1}}
  |u_0(w_0)-u_1(w_1)|>3\varepsilon,\qquad
  \inf_{\substack{w_0\in U\setminus V_0\\w_1\in U}}
  |u_0(w_0)-u_1(w_1)|>3\varepsilon.
\]
令 $v_0,v_1:U\to\mathbb C^2$ 为引理~\ref{lem:E-2-3} 所构造的浸入。
那么，除非 $w_0\in V_0$ 且 $w_1\in V_1$，否则
$v_0(w_0)\ne v_1(w_1)$。因此，两映射只在都全纯的区域相交。
由第一步，每个交点对交数 $v_0\cdot v_1$ 的贡献至少为 $1$。
此外，按构造，$v_i$ 恰有 $k_i$ 条分支经过原点。所以
\[
  \iota(u_0,u_1;0,0)=v_0\cdot v_1\ge k_0k_1>1,
\]
如所需。这证明了命题。
\end{proof}

至此完成\BookRecordIndex{T0201}局部交数的讨论。要证明\BookRecordIndex{T0055}伴随公式，还需要一个整体扰动结果。

\begin{proposition}\label{prop:E-2-4}
设 $(M,J)$ 为四维殆复流形，$(\Sigma,j)$ 为闭黎曼曲面，
$u:\Sigma\to M$ 为单纯 \BookRecordIndex{T0003}$J$-全纯曲线。取 $M$ 上使 $J$ \BookRecordIndex{T0350}斜伴随的
黎曼度量 $\langle\cdot,\cdot\rangle$，并记
\[
  \omega:=\langle J\cdot,\cdot\rangle\in\Omega^2(M).
\]
（不要求此形式闭。）则对每个 $\varepsilon>0$ 及 $\Sigma$ 中\BookRecordIndex{T0119}非单射点集
的每个邻域 $U$，存在浸入 $v:\Sigma\to M$，满足：
\begin{enumerate}
  \item[(i)] $v^*\omega$ 是 $\Sigma$ 上的\BookRecordIndex{T0404}正面积形式；
  \item[(ii)] $\|u-v\|_{C^1}<\varepsilon$；
  \item[(iii)] 对 $z\in\Sigma\setminus U$，有 $v(z)=u(z)$；
  \item[(iv)] $v$ 仅有\BookRecordIndex{T0167}横截自交点；
  \item[(v)] $v$ 满足不等式
    \begin{equation}\label{eq:E-2-3}
      \sum_{\substack{z_0\ne z_1\\v(z_0)=v(z_1)}}
        \iota(v,v;z_0,z_1)
      \ge
      \sum_{\substack{z_0\ne z_1\\u(z_0)=u(z_1)}}
        \iota(u,u;z_0,z_1),
    \end{equation}
    且等号成立当且仅当 $u$ 为浸入。
\end{enumerate}
\end{proposition}

\begin{proof}
记 $z_1,\ldots,z_N$ 为 $u$ 的奇点。由定理~\ref{thm:E-1-1}，
存在 $u(z_i)$ 的开邻域 $W_i\subset M$、$z_i$ 的开邻域
$U_i\subset\Sigma$、$C^1$-坐标图 $\Psi_i:W_i\to\mathbb C^2$、
$C^2$-坐标图 $\phi_i:U_i\to\mathbb C$、整数 $k_i>1$ 和多项式
$p_i:\mathbb C\to\mathbb C$，满足
$\phi_i(z_i)=0$、$\Psi_i(u(z_i))=0$，以及\sourcefootnote{原书此处将 $\Psi_i(u(z_i))=0$ 与 $u(U_i)\subset W_i$ 分别印作 $\Psi_i(u_i(z_i))=0$、$u_i(U_i)\subset W_i$。$u_i$ 在下文才定义为坐标表达式，不能以 $z_i$ 或 $U_i$ 为自变量；据映射定义改正。}
\[
  U_i\subset U,\qquad u(U_i)\subset W_i,\qquad
  d\Psi_i(u(z_i))J(u(z_i))=J_0d\Psi_i(u(z_i)),\qquad p_i(0)=0,
\]
并且
\[
  u_i(z):=\Psi_i\circ u\circ\phi_i^{-1}(z)
  =(z^{k_i},z^{k_i}p_i(z)),\qquad z\in\phi_i(U_i).
\]
必要时缩小 $U_i,W_i$，可设 $W_i$ 上的 $C^1$-\BookRecordIndex{T0088}殆复结构
$\Psi_i^*J_0$ 受未必闭的二形式 $\omega$ \BookRecordIndex{T0372}驯服。
对每个 $i$，固定光滑\BookRecordIndex{T0188}截断函数
$\beta_i:\phi_i(U_i)\to[0,1]$，使其在边界附近为零、在原点附近为一。
对 $\varepsilon>0$，定义 $u_\varepsilon:\Sigma\to M$ 为
\[
  u_\varepsilon(z):=\begin{cases}
    \Psi_i^{-1}\circ u_{i,\varepsilon}\circ\phi_i(z),&z\in U_i,\\
    u(z),&z\in\Sigma\setminus\bigcup_iU_i,
  \end{cases}
\]
其中
\[
  u_{i,\varepsilon}(z):=u_i(z)+\beta_i(z)
    \left(\prod_{j=0}^{k_i-1}(z+j\varepsilon)-z^{k_i}\right)(1,p_i(z)),
\]
如引理~\ref{lem:E-2-3} 的证明。令
\[
  V_i:=\{z\in U_i\mid\beta_i(\phi_i(z))=1\},\qquad
  V:=\bigcup_iV_i.
\]
于是 $u_\varepsilon|_{V_i}$ 是相对于 $\Psi_i^*J_0$ 的全纯浸入。
由于 $\Psi_i^*J_0$ 受 $\omega$ \BookRecordIndex{T0372}驯服，$u_\varepsilon^*\omega|_V$
是\BookRecordIndex{T0404}正面积形式。另一方面，$u^*\omega$ 在 $\Sigma\setminus V$ 上是\BookRecordIndex{T0404}正面积形式，
且 $u_\varepsilon$ 在 $C^1$ 拓扑下收敛于 $u$。故当 $\varepsilon>0$ 足够小时，
$u_\varepsilon^*\omega$ 是\BookRecordIndex{T0404}正面积形式。这样，我们证明了：当 $\varepsilon>0$
足够小时，$u_\varepsilon$ 是满足 (i)--(iii) 的 $C^1$-浸入。
由于 (i)、(ii) 是开条件，且 $u|_{\Sigma\setminus U}$ 是嵌入，
可在 $U$ 中扰动 $u_\varepsilon$，得到仅有横截\BookRecordIndex{T0116}二重点的光滑浸入
$v:\Sigma\to M$。因此 $v$ 满足 (i)--(iv)。为证明 (v)，先注意
\sourcefootnote{原书此处以 $v\cdot v=u_\varepsilon\cdot u_\varepsilon$ 记互异参数点处\BookRecordIndex{T0201}局部交数之和。若把 $v\cdot v$ 解释为同调类的自交数，通常还须计入\BookRecordIndex{T0117}法丛的第一陈类项；参见下文定理 E.1.6 的证明。这里保留原式，将点积仅按双点交数和理解。}
\[
  \sum_{\substack{z_0\ne z_1\\v(z_0)=v(z_1)}}
    \iota(v,v;z_0,z_1)
  =v\cdot v=u_\varepsilon\cdot u_\varepsilon
  =\sum_{\substack{z_0\ne z_1\\u_\varepsilon(z_0)=u_\varepsilon(z_1)}}
    \iota(u_\varepsilon,u_\varepsilon;z_0,z_1).
\]

$u_\varepsilon$ 的自交分两类：一类是 $z_0,z_1$ 不属于同一个 $U_i$
的点对（对应于 $u$ 的自交点）；另一类是某个 $i$ 下
$z_0,z_1\in U_i$ 的点对（对应于 $u$ 的\BookRecordIndex{T0233}临界点）。记
\begin{align*}
  S_{\varepsilon0}
  &:=\bigcup_i\{(z_0,z_1)\in U_i\times U_i\mid
      u_\varepsilon(z_0)=u_\varepsilon(z_1),\ z_0\ne z_1\},\\
  S_{\varepsilon1}
  &:=\{(z_0,z_1)\in\Sigma\times\Sigma\mid
      u_\varepsilon(z_0)=u_\varepsilon(z_1),\ z_0\ne z_1\}
      \setminus S_{\varepsilon0}.
\end{align*}
由定义~\ref{def:E-2-1}，有
\[
  \sum_{(z_0,z_1)\in S_{\varepsilon1}}
    \iota(u_\varepsilon,u_\varepsilon;z_0,z_1)
  =\sum_{\substack{z_0\ne z_1\\u(z_0)=u(z_1)}}
    \iota(u,u;z_0,z_1).
\]
此外，由 $u_\varepsilon$ 的构造，
\[
  \sum_{(z_0,z_1)\in S_{\varepsilon0}}
    \iota(u_\varepsilon,u_\varepsilon;z_0,z_1)\ge0,
\]
且等号成立当且仅当 $S_{\varepsilon0}=\varnothing$，即 $u$ 为浸入。
为说明这一点，注意若 $(z_0,z_1)\in S_{\varepsilon0}$ 是 $U_i$ 中的一对点，
则来自 $U_i$ 中点对的上述求和贡献至少为 $k_i(k_i-1)/2$。
命题得证。
\end{proof}

\begin{proof}[定理~\ref{thm:E-1-5} 的证明]
由于曲线 $u_0,u_1$ 的并是单纯的，由定理~\ref{thm:E-1-2} 可知
$u_0,u_1$ 的所有交点均孤立。因此
\[
  A_0\cdot A_1
  =\sum_{u_0(z_0)=u_1(z_1)}\iota(u_0,u_1;z_0,z_1).
\]
将命题~\ref{prop:E-2-2} 用于曲线 $u_0,u_1$ 的并，可知右边的和不小于
$\delta(u_0,u_1)$，且等号成立当且仅当所有交点横截。这证明了定理。
\end{proof}

\begin{proof}[定理~\ref{thm:E-1-6} 的证明]
设 $u:\Sigma\to M$ 为表示同调类 $A\in H_2(M;\mathbb Z)$ 的
\BookRecordIndex{T0003}$J$-全纯曲线。我们将证明
\begin{equation}\label{eq:E-2-4}
  A\cdot A-c_1(A)+\chi(\Sigma)
  \ge\sum_{\substack{z_0\ne z_1\\u(z_0)=u(z_1)}}
    \iota(u,u;z_0,z_1),
\end{equation}
且等号成立当且仅当 $u$ 为浸入。由命题~\ref{prop:E-2-2}，
右边的和不小于 $2\delta(u)$，且等号成立当且仅当所有自交点横截。
因此结论由~\eqref{eq:E-2-4} 得出。

为证明~\eqref{eq:E-2-4}，在 $M$ 上取任意使 $J$ \BookRecordIndex{T0350}斜伴随的黎曼度量
$\langle\cdot,\cdot\rangle$，记
$\omega:=\langle J\cdot,\cdot\rangle\in\Omega^2(M)$。
取满足命题~\ref{prop:E-2-4} 条件的浸入 $v:\Sigma\to M$。
由于 $v$ 是 $u$ 的扰动，故 $[v]=A$。又因 $v^*\omega$ 是 $\Sigma$
上的面积形式，存在分解
\[
  v^*TM\simeq T\Sigma\oplus\nu_v,
\]
其中 $\nu_v\to\Sigma$ 为辛向量丛；它在 $z\in\Sigma$ 上的纤维是
$T_{v(z)}M$ 中 $\operatorname{im}dv(z)$ 的\BookRecordIndex{T0351}辛补。并且
\[
  A\cdot A=\langle c_1(\nu_v),[\Sigma]\rangle
    +\sum_{\substack{z_0\ne z_1\\v(z_0)=v(z_1)}}
      \iota(v,v;z_0,z_1).
\]
由恒等式 $c_1(v^*TM)=c_1(T\Sigma)+c_1(\nu_v)$，得到
\[
  A\cdot A-c_1(A)+\chi(\Sigma)
  =\sum_{\substack{z_0\ne z_1\\v(z_0)=v(z_1)}}
    \iota(v,v;z_0,z_1).
\]
于是~\eqref{eq:E-2-4} 由命题~\ref{prop:E-2-4}(iv)--(v) 得出。
\end{proof}

\begin{remark}\label{rem:E-2-5}
设 $v:\Sigma\to M$ 是按命题~\ref{prop:E-2-4} 构造的曲线 $u$ 的扰动。
由于 $v$ 与 $u$ 同伦，和式\sourcefootnote{原书此处断言仅由互异参数点的\BookRecordIndex{T0201}局部交数之和等于同调类的自交数 $A\cdot A$。但本节定理 E.1.6 的证明在实体页 677 明确写出 $A\cdot A=\langle c_1(\nu_v),[\Sigma]\rangle+\sum\iota(v,v;z_0,z_1)$；一般还须计入\BookRecordIndex{T0117}法丛的 Euler／第一陈类项。此处保留原文断言，提醒读者核查。}
\[
  v\cdot v
  =\sum_{\substack{z_0\ne z_1\\v(z_0)=v(z_1)}}
    \iota(v,v;z_0,z_1)
\]
是同调类 $A:=[u_*(\Sigma)]$ 的自交数 $A\cdot A$。
可以在\BookRecordIndex{T0233}临界点 $z_0$ 处为 $u$ 定义\BookRecordIndex{T0201}局部交数 $\sigma(u;z_0)$，
用以计算 $z_0$ 对 $A$ 的总自交数的贡献。这样，$A\cdot A$ \BookRecordIndex{T0213}可分解为若干
贡献之和，每个\BookRecordIndex{T0119}非单射点各贡献一项。详见 McDuff~\srcite{256}。
\end{remark}

\begin{remark}\label{rem:E-2-6}
两条全纯曲线的交数可以解释为适当全纯函数的零点重数
（这是 Micallef--White~\srcite{287} 所采用的方法）。
更确切地，设多项式 $u:\mathbb C\to\mathbb C^2$ 具有形式
\[
  u(z)=(z^k,p(z)),
\]
其中 $k$ 为正整数，$p$ 被 $z^{k+1}$ 整除。定义
$f_u\in\mathbb C[x,y]$ 为
\[
  f_u(x,y):=\prod_{\lambda^k=x}(y-p(\lambda)).
\]
则 $f_u(x,y)=0$ 当且仅当 $(x,y)\in u(\mathbb C)$。
再定义 $m:=m(u)\in\mathbb N$ 为
\[
  m(u):=|\Gamma(u)|,\qquad
  \Gamma(u):=\{\lambda\in S^1\mid u(\lambda z)\equiv u(z)\}.
\]
于是 $m$ 整除 $k$。此外，$m(u)=1$ 当且仅当 $u$ 在零附近单射，
也当且仅当 $f_u$ 不可约（证明略）。由此可得多项式
$v:\mathbb C\to\mathbb C^2$，形如
$v(z)=(z^\ell,q(z))$，其中 $\ell:=k/m$，并满足
\[
  u(z)=v(z^m),\qquad f_u=f_v^m;
\]
而 $v$ 在原点附近单射，$f_v$ 不可约。现在可证明如下关于\BookRecordIndex{T0201}局部交数的
一般结论。

若 $u_0,u_1:(\mathbb D,0)\to(\mathbb C^2,0)$ 和
$f_0:(\mathbb D\times\mathbb D,0)\to(\mathbb C,0)$ 为全纯映射，
$f_0$ 在二变量全纯幂级数环中不可约，$u_0$ 为单射，
$f_0\circ u_1\not\equiv0$，并且对所有 $x,y\in\mathbb D$，
\[
  f_0(x,y)=0\quad\Longleftrightarrow\quad(x,y)\in u_0(\mathbb D),
\]
则 $(0,0)$ 为 $u_0,u_1$ 的孤立交点，且\BookRecordIndex{T0201}局部交数
\[
  \iota(u_0,u_1;0,0)=m(f_0\circ u_1)
\]
等于 $z=0$ 作为全纯函数 $f_0\circ u_1$ 零点的重数。
这里也省略证明（它使用 Hilbert \BookRecordIndex{T0236}零点定理和唯一分解性）。
这些观察给出了命题~\ref{prop:E-2-2} 的另一种证明。
\end{remark}
